\chapter{Agent Rule-Pack Architecture}
\label{ch:architecture}

\section{Repository path convention}

\begin{table}[htbp]
\centering
\footnotesize
\caption{Standard path convention for .clinerules packs.}
\label{tab:path-convention}
\begin{tabularx}{\textwidth}{@{}l X@{}}
\toprule
\textbf{Path} & \textbf{Meaning} \\
\midrule
\filepath{.clinerules} & Repository-wide or umbrella rules that are not owned by a single implementation subtree. \\
\filepath{src/<domain>/.clinerules} & Development and maintenance contract for one domain or service family. \\
\filepath{src/<domain>/.clinerules.runtime-monitor} & Runtime monitoring and alerting contract for the same domain or service family. \\
\filepath{src/<domain>/tests/test_clinerules_contracts.py} or equivalent & Validation harness that proves the rule packs are structurally complete and benchmarked. \\
\bottomrule
\end{tabularx}
\end{table}

\section{Frontmatter format definition}
\label{sec:frontmatter-format}

Every \artifact{.clinerules} pack MUST begin with a YAML frontmatter
block delimited by triple-dashes (\texttt{---}). This block defines
the mission, state, and governance linkage for the agent.

\begin{adrbox}{YAML-between-Dashes}
We mandate YAML frontmatter to ensure that metadata (version, domain,
conformance tags) is machine-readable by the validation harness without
requiring a full markdown parser.
\end{adrbox}

\section{Minimum contents of a development pack}

Every development rule pack must encode, at minimum:
\begin{enumerate}
  \item purpose and mission;
  \item source of truth;
  \item read-first files;
  \item current repo reality and known issue registry;
  \item definition of done;
  \item non-negotiable engineering rules;
  \item contract rules for schemas, payloads, and artefacts;
  \item integration rules for adjacent systems;
  \item pre-change checklist;
  \item post-change smoke tests;
  \item completion blockers.
\end{enumerate}

\section{Minimum contents of a runtime-monitor pack}

Every runtime-monitor pack must encode, at minimum:
\begin{enumerate}
  \item purpose and mission;
  \item source of truth;
  \item evidence sources;
  \item healthy versus unhealthy criteria;
  \item critical, high, medium, and informational conditions;
  \item condition codes;
  \item response rules and escalation path;
  \item operator triage checklist;
  \item minimum evidence required to close an alert.
\end{enumerate}

\section{Cross-domain inventory required by this specification}

\begin{table}[htbp]
\centering
\footnotesize
\caption{Target rule-pack inventory.}
\label{tab:inventory}
\begin{tabularx}{\textwidth}{@{}l l X@{}}
\toprule
\textbf{Area} & \textbf{Path} & \textbf{Role} \\
\midrule
Shared repository & \filepath{.clinerules} & Repo-wide documentation and standards baseline. \\
Generative UI & \filepath{src/generativeUI/.clinerules} & Deployment and monitoring standards for the UI stack. \\
Simula training & \filepath{src/training/.clinerules} & Implementation and migration contract for the Simula path. \\
Simula runtime & \filepath{src/training/.clinerules.runtime-monitor} & Monitoring and alerting contract for Simula runs. \\
Regulations development & \filepath{src/intelligence/.clinerules} & Governance-runtime implementation contract. \\
Regulations runtime & \filepath{src/intelligence/.clinerules.runtime-monitor} & Monitoring, alerting, and fail-closed contract for governance runtime. \\
\bottomrule
\end{tabularx}
\end{table}

\section{Rule-pack inheritance conflict detection}
\label{sec:conflict-detection}

The validation harness MUST include automated checks that ensure
domain-specific packs do not override critical regulatory constraints
with less restrictive values. This is the technical enforcement of the
``inherit, never weaken'' principle.

\subsection{Inheritance hierarchy}

Rule packs inherit in the following order (most general to most
specific):
\begin{enumerate}
  \item \textbf{Repository-level}: \filepath{.clinerules} (baseline)
  \item \textbf{Intelligence-level}: \filepath{src/intelligence/.clinerules}
    (regulations envelope)
  \item \textbf{Domain-level}: \filepath{src/<domain>/.clinerules}
    (domain-specific constraints)
\end{enumerate}

\subsection{Critical safety flags}

The following flags are designated as \textbf{safety-critical} and
may only be \emph{tightened} (made more restrictive) by child packs,
never \emph{loosened}:

\begin{table}[htbp]
\centering
\footnotesize
\caption{Safety-critical flags and their tightening direction.}
\label{tab:safety-flags}
\begin{tabularx}{\textwidth}{@{}l l X@{}}
\toprule
\textbf{Flag} & \textbf{Tightening direction} & \textbf{Violation example} \\
\midrule
\texttt{deny\_by\_default} & \texttt{false} $\rightarrow$ \texttt{true} only & Domain sets \texttt{true} $\rightarrow$ \texttt{false}. \\
\texttt{allow\_unlisted\_tools} & \texttt{true} $\rightarrow$ \texttt{false} only & Domain enables unlisted tools when parent denies. \\
\texttt{require\_identity\_envelope} & \texttt{false} $\rightarrow$ \texttt{true} only & Domain disables identity requirement. \\
\texttt{audit\_before\_action} & \texttt{false} $\rightarrow$ \texttt{true} only & Domain disables pre-action audit. \\
\texttt{circuit\_break\_enabled} & \texttt{false} $\rightarrow$ \texttt{true} only & Domain disables circuit breaker. \\
\texttt{max\_retries} & May only \emph{decrease} & Domain increases retry limit beyond parent. \\
\texttt{timeout\_ms} & May only \emph{decrease} & Domain increases timeout beyond parent. \\
\bottomrule
\end{tabularx}
\end{table}

\subsection{Automated conflict detection}

The validation harness SHALL:
\begin{enumerate}
  \item Load packs in inheritance order (Repository $\rightarrow$
    Intelligence $\rightarrow$ Domain).
  \item For each safety-critical flag, compare the child value against
    the parent value.
  \item If a child value is less restrictive than the parent, emit a
    \textbf{CONFLICT\_ERROR} with:
    \begin{itemize}
      \item the flag name;
      \item the parent path and value;
      \item the child path and attempted value;
      \item the allowed tightening direction.
    \end{itemize}
  \item Block CI/CD pipeline progression if any conflict is detected.
\end{enumerate}

\subsection{Conflict resolution workflow}

When a conflict is detected:
\begin{enumerate}
  \item The domain team MUST either:
    \begin{itemize}
      \item remove the conflicting override;
      \item or request a formal exception via the Controls Architecture
        Team.
    \end{itemize}
  \item Exceptions require documented risk acceptance and a time-bound
    remediation plan.
  \item All exceptions MUST be logged to
    \filepath{docs/schema/common/exceptions-registry.yaml} (validated
    against \filepath{exceptions-registry.schema.json}) with:
    \begin{itemize}
      \item exception ID;
      \item requesting domain;
      \item justification;
      \item expiry date;
      \item approving authority.
    \end{itemize}
\end{enumerate}
